% Gambler's ruin: probability of reaching the upper boundary N before
% the lower boundary 0, as a function of the fractional starting
% position k/N. The fair walk (p = q = 1/2) gives the straight line
% r_k = k/N; a slight upward drift (p > q) bows the curve upward, a
% downward drift bows it down. Curves plotted from the closed form
% r_k = (1 - (q/p)^k) / (1 - (q/p)^N) with N = 100, sampled on k.
\begin{figure}[htbp]
\centering
\begin{tikzpicture}
\begin{axis}[
  width=10.5cm, height=7.2cm,
  xlabel={starting fraction $k/N$},
  ylabel={overflow probability $r_k$},
  xmin=0, xmax=1, ymin=0, ymax=1,
  domain=0.0001:0.9999, samples=120,
  legend pos=north west,
  legend cell align=left,
  grid=both,
  grid style={enggray!18},
  tick label style={font=\small},
  label style={font=\small},
  legend style={font=\scriptsize},
]

% Fair walk: r = x (the straight line)
\addplot[very thick, engblack, domain=0:1, samples=2] {x};
\addlegendentry{$p=0.50$ (fair): $r=k/N$}

% Upward drift p = 0.52, q = 0.48, N = 100. Let a = q/p = 0.48/0.52.
% r(x) = (1 - a^{100 x}) / (1 - a^{100}). a^100 is tiny, so
% denominator ~ 1. a = 0.923076..., ln a = -0.080043.
% a^{100 x} = exp(100 x ln a) = exp(-8.0043 x).
\addplot[very thick, engblue, domain=0:1, samples=120]
  {(1 - exp(-8.0043*x)) / (1 - exp(-8.0043))};
\addlegendentry{$p=0.52$ (upward drift)}

% Downward drift p = 0.48, q = 0.52, N = 100. a = 0.52/0.48 = 1.08333,
% ln a = 0.080043. a^{100 x} = exp(8.0043 x). Denominator 1 - a^100 is
% large negative; ratio stays in [0,1].
\addplot[very thick, engred, domain=0:1, samples=120]
  {(1 - exp(8.0043*x)) / (1 - exp(8.0043))};
\addlegendentry{$p=0.48$ (downward drift)}

% Midpoint reference line
\draw[dashed, enggray] (axis cs:0.5,0) -- (axis cs:0.5,1);

\end{axis}
\end{tikzpicture}
\caption[Gambler's ruin: overflow probability]{The gambler's-ruin
overflow probability $r_k$, the chance a buffer that starts with $k$ of
$N=100$ tokens fills before it empties, as a function of the fractional
start $k/N$. For the fair walk ($p=q=\tfrac12$) the probability is the
straight line $r_k = k/N$: the outcome is decided by where the walk
starts. A four-percentage-point drift bends the curve sharply. With an
upward drift $p=0.52$ a buffer starting at the midpoint (dashed line)
overflows with probability near $0.98$; with a downward drift $p=0.48$
it overflows with probability near $0.02$. A small persistent bias
overwhelms the starting position once the boundaries are far apart,
which is the first-passage lesson behind sequential tests and
ruin-of-a-reserve problems. Source: computed by the authors from the
closed-form ruin probability of section 13.6.}
\label{fig:vol02ch13:gambler-ruin}
\end{figure}
